We define a reuse score $q_i$ for request $i$ as the fraction of its context already present in cache.
The admission rule accepts a request when
\begin{equation}
\label{eq:admit}
q_i \cdot r - c_i \ge \theta,
\end{equation}
where $r$ is the estimated recomputation cost, $c_i$ is the transfer cost, and $\theta$ is a threshold.
\label{sec:method}
A fixed margin $\Delta_{\mathrm{eff}}$ separates the two regimes.
The threshold is updated every epoch from observed tail latency.
A higher threshold admits fewer reuse requests, which protects the uplink at the cost of computation.
A lower threshold admits more reuse requests, which saves computation when the cache is warm.
The update uses a small step size so that the rule does not oscillate under bursty arrivals.
We initialize the threshold from the first epoch and let it adapt from there.
The rule never overrides the hard capacity limit of a node.
We do not yet model cache eviction, which we treat as future work.
We also leave multi-hop transfer out of scope, because it changes the cost model substantially.
